\documentclass[12pt]{article}
\usepackage[a4paper,margin=2cm]{geometry}
\usepackage{amsmath,amssymb}
\usepackage{xepersian}
\settextfont[
  Path=../../Homework/1401-2/HW2 - official solution/,
  Extension=.ttf,
  BoldFont=*Bd,
  ItalicFont=*It,
  BoldItalicFont=*BdIt
]{XB Niloofar}
\setdigitfont[
  Path=../../Homework/1401-2/HW2 - official solution/,
  Extension=.ttf,
  BoldFont=*Bd,
  ItalicFont=*It,
  BoldItalicFont=*BdIt
]{XB Niloofar}
\setlength{\parindent}{0pt}
\setlength{\parskip}{0.55em}

\begin{document}
\begin{center}
  {\Large\textbf{اصلاحیهٔ پاسخ رسمی پایان‌ترم ـ نیم‌سال ۱۴۰۱--۲}}
\end{center}

این برگه فقط موارد نادرست یا افتادهٔ پاسخ منتشرشده را اصلاح می‌کند و باید همراه فایل اصلی خوانده شود.

\section*{بخش تشریحی، سؤال ۳}
اگر
\[
  x_i=a+ih,\qquad h=\frac{b-a}{n},
\]
قاعدهٔ مرکب نقطهٔ میانی برابر است با
\[
  M_n
  =h\sum_{i=0}^{n-1}
  f\!\left(\frac{x_i+x_{i+1}}2\right).
\]
برای \(f\in C^2[a,b]\)، عددی \(\xi\in(a,b)\) وجود دارد که
\[
  \int_a^b f(x)\,dx-M_n
  =\frac{(b-a)h^2}{24}f''(\xi).
\]
در نتیجه کران قابل استفاده چنین است:
\[
  \left|\int_a^b f(x)\,dx-M_n\right|
  \leq
  \frac{(b-a)h^2}{24}
  \max_{x\in[a,b]}|f''(x)|.
\]

\section*{بخش چندگزینه‌ای، سؤال ۱۰}
برای درون‌یابی درجهٔ دو تابع
\(f(x)=\sqrt{x}\)
در گره‌های \(81,100,121\)، باقی‌مانده در \(x=90\) برابر است با
\[
 R_2(90)
 =\frac{f^{(3)}(\xi)}{3!}
 (90-81)(90-100)(90-121),
 \qquad \xi\in[81,121].
\]
چون
\[
  f^{(3)}(x)=\frac{3}{8x^{5/2}},
\]
داریم
\[
 |R_2(90)|
 \leq
 \frac{3}{8\cdot81^{5/2}}\,
 \frac{9\cdot10\cdot31}{6}
 \simeq 2.95305594\times10^{-3}.
\]
پس گزینهٔ درست \textbf{ج}، یعنی
\(2.95\times10^{-3}\)،
است.
\end{document}
